\clearpage
\VerdierPhysicalPageAnchor{0132}%
\begingroup
\small
\setlength{\emergencystretch}{1.5em}%
\setlength{\abovedisplayskip}{2pt plus 1pt minus 1pt}%
\setlength{\belowdisplayskip}{2pt plus 1pt minus 1pt}%

\VerdierSemanticAnchor{ch2-2-1-10-continuation}
\noindent 得到 SM3)。
\VerdierSemanticAnchor{ch2-2-1-10-triangle-setup}%
由假设，存在特异三角 $(X,Y,N,s,h,u)$，其中 $N$ 是 $\mathcal D'$ 的对象。
由于 $fs=0$，由 (1.2.1) 得到态射 $g\colon N\to Z$，使 $f=gh$。

\VerdierSemanticAnchor{ch2-2-1-10-factorization}
设 $(N,Z,W,g,t,v)$ 为包含 $g$ 的特异三角 (TRI)。于是 $tg=0$
(1.2.2)，从而 $tf=tgh=0$。

\VerdierSemanticAnchor{ch2-2-1-10-closure-conclusion}
所以只需证明 $t$ 属于 $S(\mathcal D')$；这立即由 TRII) 得到。

\VerdierSemanticAnchor{ch2-2-1-11}
\medskip
\noindent\textbf{2.1.11.}\quad
\VerdierSemanticAnchor{ch2-2-1-11-sm5}
\emph{集合 $S(\mathcal D')$ 具有性质} SM5)。
\VerdierSemanticAnchor{ch2-2-1-11-proof}
这立即由 TRII) 以及子范畴 $\mathcal D'$ 在平移下稳定这一事实得到。

\VerdierSemanticAnchor{ch2-2-1-12}
\medskip
\noindent\textbf{2.1.12.}\quad
\VerdierSemanticAnchor{ch2-2-1-12-sm6}
\emph{集合 $S(\mathcal D')$ 具有性质} SM6)。
\VerdierSemanticAnchor{ch2-2-1-12-triangle-pair}
设 $(X,Y,Z,u,v,w)$ 与 $(X',Y',Z',u',v',w')$ 为 $\mathcal D$ 中的两个
特异三角，并设

\VerdierSemanticAnchor{ch2-2-1-12-input-square}
\begin{center}
  \begin{tikzcd}[column sep=4.2em,row sep=2.5em]
    X \arrow[r,"u"] \arrow[d,"s"'] &
    Y \arrow[d,"{s'}"] \\
    X' \arrow[r,"{u'}"] & Y'
  \end{tikzcd}
\end{center}

\VerdierSemanticAnchor{ch2-2-1-12-completion-goal}
\noindent 为交换图，其中 $s$ 与 $s'$ 属于 $S(\mathcal D')$。需要找到
$S(\mathcal D')$ 中的态射 $s''\colon Z\to Z'$，使 $(s,s',s'')$ 是三角的态射。

\VerdierSemanticAnchor{ch2-2-1-12-case-reduction}
首先可以假设处于下列两种情形之一：

\begin{enumerate}
\renewcommand{\labelenumi}{\emph{\alph{enumi})}}
  \item \VerdierSemanticAnchor{ch2-2-1-12-case-a}%
  $X=X'$ 且 $s=\operatorname{id}_X$。
  \item \VerdierSemanticAnchor{ch2-2-1-12-case-b}%
  $Y=Y'$ 且 $s'=\operatorname{id}_Y$。
\end{enumerate}

\VerdierSemanticAnchor{ch2-2-1-12-factorization}
事实上，态射范畴中的每个态射都可作如下因式分解：

\VerdierSemanticAnchor{ch2-2-1-12-factorization-diagram}
\begin{center}
  \begin{tikzcd}[column sep=4.2em,row sep=1.7em]
    X \arrow[r,"u"] \arrow[d,"\operatorname{id}_X"'] &
    Y \arrow[d,"{s'}"] \\
    X \arrow[r,"{s'u}"] \arrow[d,"s"'] &
    Y' \arrow[d,"\operatorname{id}_{Y'}"] \\
    X' \arrow[r,"{u'}"] & Y'
  \end{tikzcd}
  \quad ,
\end{center}

\VerdierSemanticAnchor{ch2-2-1-12-composition}
\noindent 而态射 $s''$ 可通过复合得到（已经证明 $S(\mathcal D')$ 对复合
封闭）。

\VerdierSemanticAnchor{ch2-2-1-12-opposite-category}
其次，取反范畴便可由情形 a) 得到情形 b)。

\VerdierSemanticAnchor{ch2-2-1-12-triv-conclusion}
因此只需处理情形 a)，而它是 TRIV) 的直接推论。

\VerdierSemanticAnchor{ch2-2-1-13-p133-boundary}
\endgroup
